% Source-aligned English for printed p.121.
% Decisions: TOH-P121-R1-EN-001 through TOH-P121-R1-EN-009.
\SourcePage{121}

\providecommand{\TohokuContentsChapter}[1]{%
  \par\medskip\noindent\textbf{#1}\par}
\providecommand{\TohokuContentsEntry}[2]{%
  \par\noindent\hangindent=2.8em\hangafter=1%
  \makebox[2.8em][l]{#1}#2\par}

\LayerTarget{tohoku:introduction:III:p1}
of sheaves, a satisfactory account can be found in [9]. Many other gaps will
make themselves apparent to the reader.

\LayerTarget{tohoku:introduction:III:p2}
To conclude, I am happy to express my thanks to Messrs. R. Godement,
H. Cartan, and J. P. Serre, whose interest provided the indispensable stimulus
for the writing of the present work.

\LayerTarget{tohoku:contents}
\LayerContents{Table of contents}
\begin{center}
  \bfseries TABLE OF CONTENTS
\end{center}

\begingroup\small
\LayerTarget{tohoku:chapter:I}
\TohokuContentsChapter{CHAPTER I\quad GENERALITIES ON ABELIAN CATEGORIES}

\LayerTarget{tohoku:section:1.1}
\TohokuContentsEntry{1.1.}{Categories}
\LayerTarget{tohoku:section:1.2}
\TohokuContentsEntry{1.2.}{Functors}
\LayerTarget{tohoku:section:1.3}
\TohokuContentsEntry{1.3.}{Additive categories}
\LayerTarget{tohoku:section:1.4}
\TohokuContentsEntry{1.4.}{Abelian categories}
\LayerTarget{tohoku:section:1.5}
\TohokuContentsEntry{1.5.}{Infinite sums and products}
\LayerTarget{tohoku:section:1.6}
\TohokuContentsEntry{1.6.}{Categories of diagrams, permanence properties}
\LayerTarget{tohoku:section:1.7}
\TohokuContentsEntry{1.7.}{Examples of categories defined by diagram schemes}
\LayerTarget{tohoku:section:1.8}
\TohokuContentsEntry{1.8.}{Inductive and projective limits}
\LayerTarget{tohoku:section:1.9}
\TohokuContentsEntry{1.9.}{Generators and cogenerators}
\LayerTarget{tohoku:section:1.10}
\TohokuContentsEntry{1.10.}{Injective and projective objects.}
\LayerTarget{tohoku:section:1.11}
\TohokuContentsEntry{1.11.}{Quotient categories}

\LayerTarget{tohoku:chapter:II}
\TohokuContentsChapter{CHAPTER II\quad HOMOLOGICAL ALGEBRA IN ABELIAN
CATEGORIES}

\LayerTarget{tohoku:section:2.1}
\TohokuContentsEntry{2.1.}{$\partial$-functors and
$\partial^{*}$-functors}
\LayerTarget{tohoku:section:2.2}
\TohokuContentsEntry{2.2.}{Derived $\partial$-functors}
\LayerTarget{tohoku:section:2.3}
\TohokuContentsEntry{2.3.}{Derived functors}
\LayerTarget{tohoku:section:2.4}
\TohokuContentsEntry{2.4.}{Spectral sequences and spectral functors}
\LayerTarget{tohoku:section:2.5}
\TohokuContentsEntry{2.5.}{Resolving functors}

\LayerTarget{tohoku:chapter:III}
\TohokuContentsChapter{CHAPTER III\quad COHOMOLOGY WITH COEFFICIENTS IN A
SHEAF}

\LayerTarget{tohoku:section:3.1}
\TohokuContentsEntry{3.1.}{Generalities on sheaves}
\LayerTarget{tohoku:section:3.2}
\TohokuContentsEntry{3.2.}{Definition of $H_{\Phi}^{p}(X,F)$}
\LayerTarget{tohoku:section:3.3}
\TohokuContentsEntry{3.3.}{Acyclicity criteria}
\LayerTarget{tohoku:section:3.4}
\TohokuContentsEntry{3.4.}{Applications to questions of lifting the structure
group}
\LayerTarget{tohoku:section:3.5}
\TohokuContentsEntry{3.5.}{The exact sequence relative to a closed subspace}
\LayerTarget{tohoku:section:3.6}
\TohokuContentsEntry{3.6.}{On the cohomological dimension of certain spaces}
\LayerTarget{tohoku:section:3.7}
\TohokuContentsEntry{3.7.}{The Leray spectral sequence of a continuous map}
\LayerTarget{tohoku:section:3.8}
\TohokuContentsEntry{3.8.}{Comparison with Čech cohomology}
\LayerTarget{tohoku:section:3.9}
\TohokuContentsEntry{3.9.}{Acyclicity criteria by the method of coverings}
\LayerTarget{tohoku:section:3.10}
\TohokuContentsEntry{3.10.}{Passage to the limit in sheaf cohomology}

\LayerTarget{tohoku:chapter:IV}
\TohokuContentsChapter{CHAPTER IV\quad EXT FOR SHEAVES OF MODULES}

\LayerTarget{tohoku:section:4.1}
\TohokuContentsEntry{4.1.}{The functors $\mathrm{Hom}_{O}(A,B)$ and
$\mathbf{Hom}_{O}(A,B)$}
\LayerTarget{tohoku:section:4.2}
\TohokuContentsEntry{4.2.}{The functors $\mathrm{Ext}_{O}^{n}(X;\,A,B)$ and
$\mathbf{Ext}_{O}^{n}(A,B)$ and the spectral sequence}
\endgroup

\LayerTarget{tohoku:introduction:footnote:1bis}
\par\medskip\hrule\smallskip
\noindent\textsuperscript{1 bis)} M. P. Cartier has just found a satisfactory
general formulation for multiplicative structures in homological algebra,
which he will expound in due course.

\TohokuContinuation{122}
